\chapter{The Gallery of Saint-Mande}

\lettrine{F}{ifty} persons were waiting for the superintendent. He  did not even take the time to place himself in the hands of his valet de  chambre for a minute, but from the perron went straight into the premier salon.  There his friends were assembled in full chat. The intendant was about to order  supper to be served, but, above all, the Abbé Fouquet watched for the return of  his brother, and was endeavouring to do the honours of the house in his absence.  Upon the arrival of the superintendent, a murmur of joy and affection was  heard; Fouquet, full of affability, good humour, and munificence, was beloved by  his poets, his artists, and his men of business. His brow, upon which his  little court read, as upon that of a god, all the movements of his soul, and  thence drew rules of conduct,—his brow, upon which affairs of state never  impressed a wrinkle, was this evening paler than usual, and more than one  friendly eye remarked that pallor. Fouquet placed himself at the head of the  table, and presided gayly during supper. He recounted Vatel's expedition  to La Fontaine, he related the history of Menneville and the skinny fowl to  Pelisson, in such a manner that all the table heard it. A tempest of laughter  and jokes ensued, which was only checked by a serious and even sad gesture from  Pelisson. The Abbé Fouquet, not being able to comprehend why his brother should  have led the conversation in that direction, listened with all his ears, and  sought in the countenance of Gourville, or in that of his brother, an  explanation which nothing afforded him. Pelisson took up the  matter:—<Did they mention M. Colbert, then?> said he.

<Why not?> replied Fouquet; <if true, as it is said to be,  that the king has made him his intendant?> Scarcely had Fouquet uttered  these words, with a marked intention, than an explosion broke forth among the  guests.

<The miser!> said one.

<The mean, pitiful fellow!> said another.

<The hypocrite!> said a third.

Pelisson exchanged a meaning look with Fouquet. <Messieurs,> said  he, <in truth we are abusing a man whom no one knows: it is neither  charitable nor reasonable; and here is monsieur le surintendant, who, I am  sure, agrees with me.>

<Entirely,> replied Fouquet. <Let the fat fowls of M. Colbert  alone; our business to-day is with the faisans truffes of M. Vatel.> This  speech stopped the dark cloud which was beginning to throw its shade over the  guests. Gourville succeeded so well in animating the poets with the vin de  Joigny; the Abbé, intelligent as a man who stands in need of his host's  money, so enlivened the financiers and the men of the sword, that, amidst the  vapours of this joy and the noise of conversation, inquietudes disappeared  completely. The will of Cardinal Mazarin was the text of the conversation at  the second course and dessert; then Fouquet ordered bowls of sweetmeats and  fountains of liquor to be carried into the salon adjoining the gallery. He led  the way thither, conducting by the hand a lady, the queen, by his preference,  of the evening. The musicians then supped, and the promenades in the gallery  and the gardens commenced, beneath a spring sky, mild and flower-scented.  Pelisson then approached the superintendent, and said: <Something  troubles monseigneur?>

<Greatly,> replied the minister; <ask Gourville to tell you  what it is.> Pelisson, on turning round, found La Fontaine treading upon  his heels. He was obliged to listen to a Latin verse, which the poet had  composed upon Vatel. La Fontaine had, for an hour, been scanning this verse in  all corners, seeking some one to pour it out upon advantageously. He thought he  had caught Pelisson, but the latter escaped him; he turned towards Sorel, who  had, himself, just composed a quatrain in honour of the supper, and the  Amphytrion. La Fontaine in vain endeavoured to gain attention to his verses;  Sorel wanted to obtain a hearing for his quatrain. He was obliged to retreat  before M. le Comte de Charost, whose arm Fouquet had just taken. L'Abbé  Fouquet perceived that the poet, absent-minded, as usual, was about to follow  the two talkers; and he interposed. La Fontaine seized upon him, and recited  his verses. The Abbé, who was quite innocent of Latin, nodded his head, in  cadence, at every roll which La Fontaine impressed upon his body, according to  the undulations of the dactyls and spondees. While this was going on, behind  the confiture-basins, Fouquet related the event of the day to his son-in-law,  M. de Charost. <We will send the idle and useless to look at the  fireworks,> said Pelisson to Gourville, <whilst we converse  here.>

<So be it,> said Gourville, addressing four words to Vatel. The  latter then led towards the gardens the major part of the beaux, the ladies and  the chatterers, whilst the men walked in the gallery, lighted by three hundred  wax-lights, in the sight of all; the admirers of fireworks all ran away towards  the garden. Gourville approached Fouquet, and said: <Monsieur, we are  here.>

<All?> said Fouquet.

<Yes,—count.> The superintendent counted; there were eight  persons. Pelisson and Gourville walked arm in arm, as if conversing upon vague  and frivolous subjects. Sorel and two officers imitated them, and in an  opposite direction. The Abbé Fouquet walked alone. Fouquet, with M. de Charost,  walked as if entirely absorbed in the conversation of his son-in-law.  <Messieurs,> said he, <let no one of you raise his head as he  walks, or appear to pay attention to me; continue walking, we are alone, listen  to me.>

A perfect silence ensued, disturbed only by the distant cries of the joyous  guests, from the groves whence they beheld the fireworks. It was a whimsical  spectacle this, of these men walking in groups, as if each one was occupied  about something, whilst lending attention really only to one amongst them, who,  himself, seemed to be speaking only to his companion. <Messieurs,>  said Fouquet, <you have, without doubt, remarked the absence of two of my  friends this evening, who were with us on Wednesday. For God's sake,  Abbé, do not stop,—it is not necessary to enable you to listen; walk on,  carrying your head in a natural way, and as you have excellent sight, place  yourself at the window, and if any one returns towards the gallery, give us  notice by coughing.>

The Abbé obeyed.

<I have not observed their absence,> said Pelisson, who, at this  moment, was turning his back to Fouquet, and walking the other way.

<I do not see M. Lyodot,> said Sorel, <who pays me my  pension.>

<And I,> said the Abbé, at the window, <do not see M.  d'Eymeris, who owes me eleven hundred livres from our last game of  brelan.>

<Sorel,> continued Fouquet, walking bent, and gloomily, <you  will never receive your pension any more from M. Lyodot; and you, Abbé, will  never be paid you eleven hundred livres by M. d'Eymeris; for both are  doomed to die.>

<To die!> exclaimed the whole assembly, arrested, in spite of  themselves, in the comedy they were playing, by that terrible word.

<Recover yourselves, messieurs,> said Fouquet, <for perhaps  we are watched—I said: to die!>

<To die!> repeated Pelisson; <what, the men I saw six days  ago, full of health, gayety, and the spirit of the future! What then is man,  good God! that disease should thus bring him down all at once!>

<It is not a disease,> said Fouquet.

<Then there is a remedy,> said Sorel.

<No remedy. Messieurs de Lyodot and D'Eymeris are on the eve of  their last day.>

<Of what are these gentlemen dying, then?> asked an officer.

<Ask of him who kills them,> replied Fouquet.

<Who kills them? Are they being killed, then?> cried the terrified  chorus.

<They do better still; they are hanging them,> murmured Fouquet, in  a sinister voice, which sounded like a funeral knell in that rich gallery,  splendid with pictures, flowers, velvet, and gold. Involuntarily every one  stopped; the Abbé quitted his window; the first fuses of the fireworks began to  mount above the trees. A prolonged cry from the gardens attracted the  superintendent to enjoy the spectacle. He drew near to a window, and his  friends placed themselves behind him, attentive to his least wish.

<Messieurs,> said he, <M. Colbert has caused to be arrested,  tried and will execute my two friends; what does it become me to do?>

<Mordieu!> exclaimed the Abbé, the first one to speak, <run  M. Colbert through the body.>

<Monseigneur,> said Pelisson, <you must speak to his  majesty.>

<The king, my dear Pelisson, himself signed the order for the  execution.>

<Well!> said the Comte de Charost, <the execution must not  take place, then; that is all.>

<Impossible,> said Gourville, <unless we could corrupt the  jailers.>

<Or the governor,> said Fouquet.

<This night the prisoners might be allowed to escape.>

<Which of you will take charge of the transaction?>

<I,> said the Abbé, <will carry the money.>

<And I,> said Pelisson, <will be the bearer of the  words.>

<Words and money,> said Fouquet, <five hundred thousand  livres to the governor of the conciergerie that is sufficient; nevertheless, it  shall be a million, if necessary.>

<A million!> cried the Abbé; <why, for less than half, I  would have half Paris sacked.>

<There must be no disorder,> said Pelisson. <The governor  being gained, the two prisoners escape; once clear of the fangs of the law,  they will call together the enemies of Colbert, and prove to the king that his  young justice, like all other monstrosities, is not infallible.>

<Go to Paris, then, Pelisson,> said Fouquet, <and bring  hither the two victims; to-morrow we shall see.>

Gourville gave Pelisson the five hundred thousand livres. <Take care the  wind does not carry you away,> said the Abbé; <what a  responsibility. Peste! Let me help you a little.>

<Silence!> said Fouquet, <somebody is coming. Ah! the  fireworks are producing a magical effect.> At this moment a shower of  sparks fell rustling among the branches of the neighbouring trees. Pelisson and  Gourville went out together by the door of the gallery; Fouquet descended to  the garden with the five last plotters.  